\section{Multilinearity in tangent categories}
\label{section:linearity}
To introduce general linear group objects in tangent categories, we first take a detour and study internal multilinear Homs.



%__________________________________________________________________________
\subsection{A fibrational perspective on multilinearity}
\label{subsection:fibrational-multilinearity}
In~\cite{lemay:jacobian}, Lemay introduced \emph{linearly closed Cartesian differential categories} to give a synthetic perspective on the Jacobian matrix of smooth morphisms. In this section, we extend this concept to Cartesian tangent categories. We begin by adapting to our context the notions of linear and multilinear morphisms, introduced by Cruttwell in~\cite[Definition~2.9]{cruttwell:sector-forms-in-CDCs} in the context of Cartesian differential categories.

\begin{definition}[{\cite[Definition~2.9, adapted for tangent categories]{cruttwell:sector-forms-in-CDCs}}]
\label{definition:multilinearity}
In a Cartesian tangent category:
\begin{enumerate}
\item A morphism $f\colon M_1\ntimes\A_k\ntimes M_n\to\B$ is \textbf{linear in the $k$-th argument}, where $\A_k$ and $\B$ are differential objects, provided that
\begin{equation*}
% https://q.uiver.app/#q=WzAsNCxbMCwwLCJcXFQoTV8xXFx0aW1lc1xcZG90c1xcdGltZXNcXEFfa1xcdGltZXNcXGRvdHNcXHRpbWVzIE1fbikiXSxbMCwxLCJNXzFcXHRpbWVzXFxkb3RzXFx0aW1lc1xcQV9rXFx0aW1lc1xcZG90c1xcdGltZXMgTV9uIl0sWzEsMSwiXFxCIl0sWzEsMCwiXFxUXFxCIl0sWzEsMCwiXFxsYW1iZGFfayJdLFsyLDMsIlxcbGFtYmRhX1xcQiIsMl0sWzEsMiwiZiIsMl0sWzAsMywiXFxUIGYiXV0=
\begin{tikzcd}
{\T(M_1\ntimes\A_k\ntimes M_n)} & {\T\B} \\
{M_1\ntimes\A_k\ntimes M_n} & \B
\arrow["{\T f}", from=1-1, to=1-2]
\arrow["{\lambda_k}", from=2-1, to=1-1]
\arrow["f"', from=2-1, to=2-2]
\arrow["{\lambda}"', from=2-2, to=1-2]
\end{tikzcd}
\end{equation*}
where $\lambda_k\=z_{M_1}\ntimes\lambda_{\A_k}\ntimes z_{M_n}$.

\item A morphism $f\colon\A_1\ntimes\A_n\to\B$ is $(\A_1\,\A_n)$-\textbf{multilinear} (or simply \textbf{multilinear}), where $\A_1\,\A_n,\B$ are differential objects, provided that $f$ is linear in each $k$-th argument, for $k=1\,n$.

\item An $(\A)$-multilinear morphism $f\colon\A\to\B$ is called \textbf{linear}\footnote{This corresponds to the usual definition of linearity between differential objects introduced in~\cite{cockett:differential-bundles}.}.

\item An $(\A_1,\A_2)$-multilinear morphism $f\colon\A_1\times\A_2\to\B$ is called \textbf{bilinear}.
\end{enumerate}
\end{definition}

Differential objects together with multilinear maps form a multicategory denoted by $\DO_\multi[\X,\TT]$.

We shall spend some time reinterpreting multilinearity through a fibrational lens. Our goal is to express multilinearity as proper \emph{linearity at each entry}. We begin by recalling that the \textbf{simple fibration} of a Cartesian category $\X$ (introduced in~\cite[Section~1.3]{jacobs:fibrations-categorical-logic-type-theory}) is the (cloven) fibration
\begin{align*}
&\Pi\colon\Simple[\X]\to\X
\end{align*}
whose total category $\Simple[\X]$ is the category of pairs $(M,A)$ of objects of $\X$, whose morphisms are pairs $(f,\varphi)$ of morphisms $f\colon M\to N$ and $\varphi\colon A\times M\to B$\footnote{We adopted the non-standard convention of switching between $A$ and $M$. This convention is more natural when we consider left-invariant vector fields.}. The identity morphism of $(M,A)$ is given by $(\id_M,\pi_1)$ and composition is defined by the following formula:
\begin{align*}
&(f,\varphi)(g,\psi)=\left(fg,A\times M\xrightarrow{\id_A\times\Delta_M}A\times M\times M\xrightarrow{\varphi\times f}B\times N\xrightarrow{\psi}C\right)
\end{align*}
The fibration $\Pi$ sends each $(M,A)$ to $M$ and each $(f,\varphi)$ to $f$. When the Cartesian category $\X$ is equipped with a Cartesian tangent structure $\TT$, $\Pi$ becomes a (cloven) \textbf{tangent fibration}~\cite[Definition~5.2]{cockett:differential-bundles}. Recall that a (cloven) tangent fibration consists of a fibration between two tangent categories whose underlying functor preserves the tangent structure strictly and the tangent bundle functors preserve the (cloven) Cartesian lifts.

\begin{proposition}[{\cite[Proposition~3.7]{lanfranchi:lie-correspondence-tangent-cats}}]
\label{proposition:simple-tangent-fibration}
The simple fibration $\Pi\colon\Simple[\X,\TT]\to(\X,\TT)$ of a Cartesian tangent category is a cloven tangent fibration.
\end{proposition}

By~\cite[Theorem~5.3]{cockett:differential-bundles}, the fibres $\Simple_M[\X,\TT]$ become tangent categories and the substitution functors between the fibres become strong tangent morphisms. Concretely, the tangent bundle functor of $\Simple[\X,\TT]$ sends each $(M,A)$ to $(\T M,\T A)$ and each $(f,\varphi)$ to $(\T f,\varphi_\T)$, where:
\begin{align*}
&\varphi_\T\colon\T A\times\T M\xrightarrow{\cong}\T(A\times M)\xrightarrow{\T\varphi}\T B
\end{align*}
Furthermore, for each $M\in(\X,\TT)$, the fibre $\Simple_M[\X,\TT]$ over $M$ of the simple tangent fibration becomes a tangent category whose objects are the objects of $\X$, morphisms are morphisms $\varphi\colon A\times M\to B$, and the tangent bundle functor sends each $A$ to $\T A$ and each $\varphi$ to:
\begin{align*}
&\T A\times M\xrightarrow{\id_{\T A}\times z_M}\T(A\times M)\xrightarrow{\T\varphi}\T B
\end{align*}

\begin{lemma}[{\cite[Lemma~4.13]{lanfranchi:lie-correspondence-tangent-cats}}]
\label{lemma:inclusion-functor-preserves-differential-objects}
The inclusion functor $(\X,\TT)\to\Simple_M[\X,\TT]$ that sends each $A$ to itself and each $f\colon A\to B$ to $f^\KL$ defined by $\SQ{f^\KL}\=\pi_1f\colon A\times M\to B$ becomes a strict tangent morphism that preserves tangent limits. Moreover, it preserves differential objects.
\end{lemma}

We shall denote by $\A_1\ntimes\hat\A_k\ntimes\A_n$ the Cartesian product of a list of differential objects in which the $k$-th term has been omitted.

\begin{lemma}
\label{lemma:multilinearity-fibrationally}
A morphism $f\colon M_1\ntimes\A_k\ntimes M_n\to\B$ is linear in the $k$-th argument if and only if
\begin{align*}
\SQ{f_k}\colon\A_k\times M_1\ntimes\hat\A_k\ntimes M_n\xrightarrow{\Shuffle_k}M_1\ntimes\A_k\ntimes M_n\xrightarrow{f}\B
\end{align*}
lifts to a linear morphism $\A_k\nto\B$ in $\Simple_{M_1\ntimes\hat\A_k\ntimes M_n}[\X,\TT]$, where $\Shuffle_k$ denotes the permutation that sends the first entry to the $k$-th entry and preserves the orders of the other entries.
\end{lemma}
\begin{proof}
Consider a morphism $f\colon M_1\ntimes\A_k\ntimes M_n\to\B$ and suppose that $f_k\colon\A_k\nto\B$ is linear in $\Simple_k[\X,\TT]\=\Simple_{M_1\ntimes\hat\A_k\ntimes M_n}[\X,\TT]$. By definition, this means that $f_k\lambda_\B=\lambda_{\A_k}\T_{(k)}f_k$ in $\Simple_k[\X,\TT]$. By construction, the tangent bundle functor $\T_{(k)}$ sends $f_k$ to:
\begin{align*}
&\SQ{\T_{(k)}f_k}=(\id\times z\ntimes z)\T f_k=(\id\times z\ntimes z)\Shuffle_k\T f=\\
&\qquad=\Shuffle_k(z\ntimes\id\ntimes z)\T f
\end{align*}
Thus, using the formula that defines the coKleisli composition in $\Simple_k[\X,\TT]$, we rewrite the linearity condition as follows:
\begin{align*}
&\Shuffle_k f\lambda_\B=\Shuffle_k(z\ntimes\lambda_{\A_k}\ntimes z)\T f=\Shuffle_k\lambda_k\T f
\end{align*}
However, since $\Shuffle_k$ is just a permutation, it is invertible, thus $fl\lambda_\B=\lambda_k\T f$, which means that $f$ is linear in the $k$-th argument.
\end{proof}

As explained in the introduction, a differential object is a commutative monoid $\A$ equipped with a vertical lift $\lambda\colon\A\to\T\A$ that is used to construct an isomorphism $\xi\colon\A\times\A\to\T\A$. Equivalently, a differential object is a commutative monoid $\A$ equipped with a differential projection $\hat p$ that turns $\T\A$ into the product $\A\times\A$. A morphism $f\colon\A\to\B$ is linear if $f\lambda_\B=\lambda_\A\T f$, or as expressed via the differential projections, if $\T f\hat p_\B=\hat p_\A f$. Thanks to Lemma~\ref{lemma:multilinearity-fibrationally}, we can rephrase multilinearity as a compatibility condition with the differential projections.

\begin{lemma}
\label{lemma:multilinearity-as-p-hat}
A morphism $f\colon M_1\ntimes\A_k\ntimes M_n\to\B$ is linear in its $k$-th argument if and only if the following diagram
\begin{equation*}
% https://q.uiver.app/#q=WzAsNSxbMCwxLCJNXzFcXG50aW1lc1xcQV9rXFxudGltZXMgTV9uIl0sWzAsMCwiTV8xXFxudGltZXNcXFRcXEFfa1xcbnRpbWVzIE1fbiJdLFszLDEsIlxcQiJdLFsyLDAsIlxcVChNXzFcXG50aW1lc1xcQV9rXFxudGltZXMgTV9uKSJdLFszLDAsIlxcVFxcQiJdLFsxLDAsIlxcaGF0IHBfayIsMl0sWzAsMiwiZiIsMl0sWzEsMywielxcbnRpbWVzIHpcXHRpbWVzXFxpZFxcdGltZXMgelxcbnRpbWVzIHoiXSxbMyw0LCJcXFQgZiJdLFs0LDIsIlxcaGF0IHBfXFxCIl1d
\begin{tikzcd}
{M_1\ntimes\T\A_k\ntimes M_n} && {\T(M_1\ntimes\A_k\ntimes M_n)} & {\T\B} \\
{M_1\ntimes\A_k\ntimes M_n} &&& \B
\arrow["{z\ntimes z\times\id\times z\ntimes z}", from=1-1, to=1-3]
\arrow["{\hat p_k}"', from=1-1, to=2-1]
\arrow["{\T f}", from=1-3, to=1-4]
\arrow["{\hat p_\B}", from=1-4, to=2-4]
\arrow["f"', from=2-1, to=2-4]
\end{tikzcd}
\end{equation*}
commutes, where:
\begin{align*}
&\hat p_k\=\id\ntimes\id\times\hat p_{\A_k}\times\id\ntimes\id
\end{align*}
In particular, a morphism $f\colon\A_1\ntimes\A_n\to\B$ is multilinear if and only if, for each $k=1\,n$ the following diagram
\begin{equation*}
% https://q.uiver.app/#q=WzAsNSxbMCwxLCJcXEFfMVxcbnRpbWVzXFxBX24iXSxbMCwwLCJcXEFfMVxcbnRpbWVzXFxUXFxBX2tcXG50aW1lcyBcXEFfbiJdLFszLDEsIlxcQiJdLFsyLDAsIlxcVChcXEFfMVxcbnRpbWVzXFxBX24pIl0sWzMsMCwiXFxUXFxCIl0sWzEsMCwiXFxoYXQgcF9rIiwyXSxbMCwyLCJmIiwyXSxbMSwzLCJ6XFxudGltZXMgelxcdGltZXNcXGlkXFx0aW1lcyB6XFxudGltZXMgeiJdLFszLDQsIlxcVCBmIl0sWzQsMiwiXFxoYXQgcF9cXEIiXV0=
\begin{tikzcd}
{\A_1\ntimes\T\A_k\ntimes \A_n} && {\T(\A_1\ntimes\A_n)} & {\T\B} \\
{\A_1\ntimes\A_n} &&& \B
\arrow["{z\ntimes z\times\id\times z\ntimes z}", from=1-1, to=1-3]
\arrow["{\hat p_k}"', from=1-1, to=2-1]
\arrow["{\T f}", from=1-3, to=1-4]
\arrow["{\hat p_\B}", from=1-4, to=2-4]
\arrow["f"', from=2-1, to=2-4]
\end{tikzcd}
\end{equation*}
commutes.
\end{lemma}
\begin{proof}
A morphism $f\colon\A\to\B$ is linear if and only if it satisfies $\T f\hat p_\B=\hat p_\A f$. From Lemma~\ref{lemma:multilinearity-fibrationally}, a morphism $f\colon M_1\ntimes\A_k\ntimes M_n\to\B$ is linear in its $k$-th entry if and only if $f_k\colon\A_k\nto\B$, defined by $\SQ{f_k}\=\Shuffle_kf$, becomes a linear morphism in $\Simple_k[\X,\TT]\=\Simple_{M_1\ntimes\hat\A_k\ntimes M_n}[\X,\TT]$. However, $f_k$ is linear if and only if $\T_{(k)}f_k\hat p_\B=\hat p_{\A_k}f_k$ in $\Simple_k[\X,\TT]$, where $\T_{(k)}$ denotes its tangent bundle functor. However, by unwrapping the definition of $\T_{(k)}$ and of the composition, we rewrite it as follows:
\begin{align*}
&\SQ{\T_{(k)}f_k\hat p_\B}=(\id\times z\ntimes z)\T f_k\hat p_\B        &&\SQ{\hat p_\A f_k}=(\hat p_{\A_k}\times\id\ntimes\id)f_k
\end{align*}
Therefore, by precomposing by $\Shuffle_k^{-1}$, we obtain the desired equation:
\begin{align*}
&(z\ntimes\id\ntimes z)\T f\hat p_\B=(\id\ntimes\hat p_{\A_k}\ntimes\id)f=\hat p_kf\qedhere
\end{align*}
\end{proof}

In the next section, we will extensively consider morphisms $f\colon\A_1\ntimes\A_n\times M\to\B$ which are linear in each $k$-th entry, for $k=1\,n$. Thanks to the fibrational perspective we presented, we can reinterpret such maps as \emph{multilinear} morphisms $f\colon\A_1\ntimes\A_n\nto\B$ in the tangent category $\Simple_M[\X,\TT]$.

\begin{definition}
\label{definition:M-parametrized-multilinear-map}
An \textbf{$M$-parametrized multilinear map} is a map $f\colon\A_1\ntimes\A_n\times M\to\B$ that lifts to a multilinear morphism $f\colon\A_1\,\A_n\to\B$ in $\Simple_M[\X,\TT]$.
\end{definition}



%__________________________________________________________________________
\subsection{Multilinearly closed tangent categories}
\label{subsection:internal-linear-homs}
We shall now introduce internal multilinear Homs. We begin by adapting Lemay's definition of internal \emph{linear} Homs in the context of tangent categories and extend it to capture multilinear maps.

\begin{definition}[{\cite[Definition~3.1, adapted for tangent categories]{lemay:jacobian}}]
\label{definition:linearly-closed-tangent-category}
In a Cartesian tangent category, given a list of differential objects $\A_1\,\A_n,\B$, the \textbf{internal multilinear Hom} from $\A_1\,\A_n$ to $\B$ consists of an object $\lnr[\A_1\,\A_n;\B]$ equipped with a morphism
\begin{align*}
&\ev\colon\A_1\ntimes\A_n\times\lnr[\A_1\,\A_n;\B]\to\B
\end{align*}
called the \textbf{evaluation map}, subject to the following conditions:
\begin{description}
\item[ILH.1] $\lnr[\A_1\,\A_n;\B]$ is a differential object and the evaluation map is multilinear, in that:
\begin{equation*}
% https://q.uiver.app/#q=WzAsNCxbMCwxLCJcXEFfMVxcbnRpbWVzXFxBX25cXHRpbWVzXFxsbnJbXFxBXzFcXCxcXEFfbjtcXEJdIl0sWzEsMSwiXFxCIl0sWzAsMCwiXFxUKFxcQV8xXFxudGltZXNcXEFfblxcdGltZXNcXGxucltcXEFfMVxcLFxcQV9uO1xcQl0pIl0sWzEsMCwiXFxUXFxCIl0sWzAsMiwiXFxsYW1iZGFfayJdLFsxLDMsIlxcbGFtYmRhX1xcQiIsMl0sWzAsMSwiXFxldl97XFxBXzFcXCxcXEFfbjtcXEJ9IiwyXSxbMiwzLCJcXFRcXGV2X3tcXEFfMVxcLFxcQV9uO1xcQn0iXV0=
\begin{tikzcd}
{\T(\A_1\ntimes\A_n\times\lnr[\A_1\,\A_n;\B])} & {\T\B} \\
{\A_1\ntimes\A_n\times\lnr[\A_1\,\A_n;\B]} & \B
\arrow["{\T\ev}", from=1-1, to=1-2]
\arrow["{\lambda_k}", from=2-1, to=1-1]
\arrow["{\ev}"', from=2-1, to=2-2]
\arrow["{\lambda_\B}"', from=2-2, to=1-2]
\end{tikzcd}
\end{equation*}
for $k=1\,n+1$;

\item[ILH.2] The evaluation map classifies multilinear maps, in that for every $M$-parametrized multilinear map $f\colon\A_1\ntimes\A_n\times M\to\B$, there exists a unique morphism $\hat f\colon M\to\lnr[\A_1\,\A_n;\B]$ that makes the following diagram
\begin{equation*}
% https://q.uiver.app/#q=WzAsMyxbMCwwLCJcXEFfMVxcbnRpbWVzXFxBX25cXHRpbWVzXFxsbnJbXFxBXzFcXCxcXEFfbjtcXEJdIl0sWzEsMCwiXFxCIl0sWzAsMSwiXFxBXzFcXG50aW1lc1xcQV9uXFx0aW1lcyBNIl0sWzAsMSwiXFxldl97XFxBXzFcXCxcXEFfbjtcXEJ9Il0sWzIsMSwiZiIsMl0sWzIsMCwiXFxpZF97XFxBXzFcXG50aW1lc1xcQV9ufVxcdGltZXNcXGhhdCBmIl1d
\begin{tikzcd}
{\A_1\ntimes\A_n\times\lnr[\A_1\,\A_n;\B]} & \B \\
{\A_1\ntimes\A_n\times M}
\arrow["{\ev}", from=1-1, to=1-2]
\arrow["{\id\times\hat f}", from=2-1, to=1-1]
\arrow["f"', from=2-1, to=1-2]
\end{tikzcd}
\end{equation*}
commutative.
\end{description}
When $n=1$, we refer to $\lnr[\A;\B]$ as the \textbf{internal linear Hom}. A tangent category is \textbf{multilinearly closed} provided that the internal multilinear Hom $\lnr[\A_1\,\A_n;\B]$ exists for each list $\A_1\,\A_n,\B$ of differential objects.
\end{definition}

\begin{remark}
\label{remark:internal-linear-homs}
Our Definition~\ref{definition:linearly-closed-tangent-category} and Lemay's original \cite[Definition~3.1]{lemay:jacobian} differ in a few details. First, the ambient category is not simply a Cartesian differential category but, in fact, a Cartesian tangent category. Second, the universal property \textbf{[ILH.2]}. In fact, this universal property classifies not only multilinear maps but also every \emph{parametrized} multilinear map, where the parameter object $M$ can be chosen arbitrarily in the tangent category. Finally, the original definition only defines internal \emph{linear} Homs, while our version extends this to general multilinear maps.
\end{remark}

Despite these differences between ours and Lemay's original definition, they are equivalent in the category of differential objects. We shall omit the easy proof of this statement.

\begin{proposition}
\label{proposition:internal-linear-hom-js}
The Cartesian differential category $\DO[\X,\TT]$ of differential objects of a multilinearly closed tangent category $(\X,\TT)$ is linearly closed in the sense of~\cite[Definition~3.1]{lemay:jacobian}. In particular, each $\lnr[\A,\B]$ becomes an internal linear Hom as in~\cite[Definition~3.1]{lemay:jacobian}.
\end{proposition}

We now show that each $\lnr[\A_1\,\A_n;\B]$ is an exponential object in the multicategory $\DO_\multi[\X,\TT]$ of differential objects and multilinear morphisms. Recall that an exponential object in a multicategory consists of an object $\B^{\A_1\,\A_n}$ equipped with a multimorphism $\ev\colon\A_1\,\A_n,\B^{\A_1\,\A_n}\to\B$ subject to the following universal property. For every multimorphism $f\colon\A_1\,\A_n,\A'_1\,\A'_m\to\B$, there exists a unique morphism $\hat f\colon\A_1'\,\A_m'\to\B^{\A_1\,\A_n}$ such that $(\id_{\A_1}\,\id_{\A_n},\hat f)\ev=f$.

\begin{proposition}
\label{proposition:exponential-objects}
Internal multilinear Homs are exponential objects in the multicategory $\DO_\multi[\X,\TT]$.
\end{proposition}
\begin{proof}
By \textbf{[ILH.1]}, each internal multilinear Hom is an object of $\DO_\multi[\X,\TT]$ and the evaluation map $\ev$ becomes a multimorphism $\ev\colon\A_1\,\A_n,\lnr[\A_1\,\A_n;\B]\to\B$ of $\DO_\multi[\X,\TT]$. It is left only to prove that for each $f\colon\A_1\,\A_n,\A_1'\,\A_m'\to\B$, $\hat f\colon\A_1'\ntimes\A_m'\to\lnr[\A_1\,\A_n;\B]$, which is the unique morphism satisfying $(\id\times\hat f)\ev=f$, is multilinear. To this end, we shall compute:
\begin{align*}
(\id\ntimes\id\times(\hat p_k\hat f))\ev&=~(\id\ntimes\id\times\hat p_k)(\id\ntimes\id\times\hat f)\ev\\
&=~(\id\ntimes\id\times\hat p_k)f                   \tag{Definition of $\hat f$}\\
&=~(z\ntimes z\times\id)\T f\hat p                  \tag{$f$ is multilinear}\\
&=~(z\ntimes z\times\id)(\id\ntimes\id\times\T\hat f)\T\ev\hat p    \tag{Definition of $\hat f$}\\
&=~(\id\ntimes\id\times\T\hat f)(z\ntimes z\times\id)\T\ev\hat p\\
&=~(\id\ntimes\id\times\T\hat f)(\id\ntimes\id\times\hat p)\ev          \tag{$\ev$ is multilinear}\\
&=~(\id\ntimes\id\times(\T\hat f\hat p))\ev
\end{align*}
By invoking \textbf{[ILH.2]}, we conclude that $\hat p_k\hat f=\T\hat f\hat p$, proving, thanks to Lemma~\ref{lemma:multilinearity-as-p-hat}, that $\hat f$ is in fact multilinear.
\end{proof}

In Definition~\ref{definition:linearly-closed-tangent-category}, \textbf{[ILH.1]} establishes that the object $\lnr[\A_1\,\A_n;\B]$ should be a differential object and that the evaluation map should be multilinear, while \textbf{[ILH.2]} establishes the class of morphisms that internal multilinear Homs classify. It turns out that one can \emph{derive} that $\lnr[\A_1\,\A_n;\B]$ is a differential object by enlarging the class of morphisms classified by the evaluation map.

\begin{theorem}
\label{theorem:internal-linear-homs-classify-kleisli-multilinear-morphisms}
Given a list $\A_1\,\A_n,\B$ of differential object, consider an object $\lnr[\A_1\,\A_n;\B]$ equipped with a map $\ev\colon\A_1\ntimes\A_n\times\lnr[\A_1\,\A_n;\B]\to\B$. Then $\lnr[\A_1\,\A_n;\B]$ is an internal multilinear Hom with evaluation map $\ev$ if and only if the following conditions are satisfied:
\begin{description}
\item[ILH.1'] The map $\ev$ is a $\lnr[\A_1\,\A_n;\B]$-parametrized multilinear map;

\item[ILH.2'] For every $M$-parametrized multilinear map $f\colon\A_1\ntimes\A_n\times M\to\T\B$, there exists a unique morphism $\hat f\colon M\to\T\lnr[\A_1\,\A_n;\B]$ that makes the following diagram
\begin{equation*}
% https://q.uiver.app/#q=WzAsMyxbMCwwLCJcXFQoXFxBXzFcXG50aW1lc1xcQV9uXFx0aW1lc1xcbG5yW1xcQV8xXFwsXFxBX247XFxCXSkiXSxbMSwwLCJcXFRcXEIiXSxbMCwxLCJcXEFfMVxcbnRpbWVzXFxBX25cXHRpbWVzIE0iXSxbMCwxLCJcXFRcXGV2Il0sWzIsMSwiZiIsMl0sWzIsMCwielxcbnRpbWVzIHpcXHRpbWVzXFxoYXQgZiJdXQ==
\begin{tikzcd}
{\T(\A_1\ntimes\A_n\times\lnr[\A_1\,\A_n;\B])} & {\T\B} \\
{\A_1\ntimes\A_n\times M}
\arrow["{\T\ev}", from=1-1, to=1-2]
\arrow["{z\ntimes z\times\hat f}", from=2-1, to=1-1]
\arrow["f"', from=2-1, to=1-2]
\end{tikzcd}
\end{equation*}
commutative.
\end{description}
\end{theorem}
\begin{proof}
We begin by assuming that $\lnr[\A_1\,\A_n;\B]$ is an internal multilinear Hom object and we prove that it satisfies \textbf{[ILH.1']} and \textbf{[ILH.2']}. \textbf{[ILH.1']} follows immediately from \textbf{[ILH.1]}. To prove \textbf{[ILH.2']}, consider a parametrized multilinear map $f\colon\A_1\ntimes\A_n\times M\to\T\B$. Since $\B$ is a differential object, $\T\B\cong\B\times\B$, so we can consider $f_1\=fp_\B$ and $f_2\=f\hat p_\B$, which are also parametrized multilinear maps. Therefore, by \textbf{[ILH.2]}, we find two morphisms $\hat f_1,\hat f_2\colon M\to\lnr[\A_1\,\A_n;\B]$, satisfying:
\begin{align*}
&(\id\times\hat f_1)\ev=f_1             &&(\id\times\hat f_2)\ev=f_2
\end{align*}
However, by \textbf{[ILH.1]}, $\T\lnr[\A_1\,\A_n;\B]\cong\lnr[\A_1\,\A_n;\B]\times\lnr[\A_1\,\A_n;\B]$. So, we can define $\hat f\colon M\to\T\lnr[\A_1\,\A_n;\B]$ as the morphism $\hat f=\<\hat f_1,\hat f_2\>$. We shall prove that $(z\ntimes z\times\hat f)\T\ev=f$:
\begin{align*}
(z\ntimes z\times\hat f)\T\ev p&=~(zp\ntimes zp\times\hat fp)\ev            \tag{Naturality of $p$}\\
&=~(\id\times\hat f_1)\ev                                                   \tag{$zp=\id$, $\hat fp=\hat f_1$}\\
&=~f_1                                                                      \tag{Definition of $\hat f_1$}
\end{align*}
Moreover:
\begin{align*}
(z\ntimes z\times\hat f)\T\ev\hat p&=~(\id\times\hat f\hat p)\ev            \tag{$\ev$ is multilinear}\\
&=~(\id\times\hat f_2)\ev                                                   \tag{$\hat f\hat p=\hat f_2$}\\
&=~f_2                                                                      \tag{Definition of $\hat f_2$}
\end{align*}
This proves that $(z\ntimes z\times\hat f)\T\ev=f$. Now, consider a morphism $g\colon M\to\T\lnr[\A_1\,\A_n;\B]$ such that $(z\ntimes z\times g)\T\ev=f$. Thus:
\begin{align*}
&(\id\times\hat f_1)\ev=f_1=fp=(z\ntimes z\times g)\T\ev p=(\id\times gp)\ev\\
&(\id\times\hat f_2)\ev=f_2=f\hat p=(z\ntimes z\times g)\T\ev\hat p=(\id\times g\hat p)\ev
\end{align*}
Therefore, by \textbf{[ILH.2]}, $\hat f_1=gp$ and $\hat f_2=g\hat p$; thus, $g=\hat f$. This proves \textbf{[ILH.2']}.

Conversely, we now assume that $\lnr[\A_1\,\A_n;\B]$ with $\ev$ satisfies \textbf{[ILH.1']} and \textbf{[ILH.2']} and we shall prove that $\lnr[\A_1\,\A_n;\B]$ is an internal multilinear Hom object. We begin by proving that \textbf{[ILH.2]} holds. To this end, consider an $M$-parametrized multilinear map $f\colon\A_1\ntimes\A_n\times M\to\B$. It is easy to see that $fz$ is also parametrized multilinear. Thus, by \textbf{[ILH.2']}, there exists a unique morphism $g\colon M\to\T\lnr[\A_1\,\A_n;\B]$ such that $(z\ntimes z\times g)\T\ev=fz$. We shall prove that $gp$ satisfies $(\id\times gp)\ev=f$:
\begin{align*}
f&=~fzp                                         \tag{$zp=\id$}\\
&=~(z\ntimes z\times g)\T\ev p                  \tag{Definition of $g$}\\
&=~(\id\times gp)\ev                            \tag{Naturality of $p$}
\end{align*}        
Now, consider another morphism $h\colon M\to\lnr[\A_1\,\A_n;\B]$ that satisfies $(\id\times h)\ev=f$. Thus:
\begin{align*}
&(z\ntimes z\times g)\T\ev=fz=(\id\times h)\ev z=(z\ntimes z\times hz)\T\ev
\end{align*}
where we used the naturality of $z$. Thus, by \textbf{[ILH.2']}, $hz=g$; $h=gp$, proving uniqueness. This shows that \textbf{[ILH.2]} holds. Thus, in the following, we shall assume that \textbf{[ILH.2]} is given. Next, we show that $\lnr[\A_1\,\A_n;\B]$ is a differential object and that $\ev$ is multilinear. We begin by constructing the additive structure of $\lnr[\A_1\,\A_n;\B]$. We define the zero $0\colon\*\to\lnr[\A_1\,\A_n;\B]$ as the unique morphism that satisfies the following equation: $(\id\times0)\ev=!0$. It is not hard to see that $!0$ is a parametrized multilinear map. Thus, the existence of $0$ follows from the universal property \textbf{[ILH.2]}. Similarly, we define the sum $+\colon\lnr[\A_1\,\A_n;\B]\times\lnr[\A_1\,\A_n;\B]\to\lnr[\A_1\,\A_n;\B]$ as the unique morphism that makes the following diagram
\begin{equation*}
\adjustbox{width=\linewidth}{
% https://q.uiver.app/#q=WzAsNSxbMiwwLCJcXEIiXSxbMCwwLCJcXEFcXHRpbWVzXFxsbnJbXFxBXzFcXCxcXEFfbjtcXEJdIl0sWzAsMiwiXFxBXzFcXG50aW1lc1xcQV9uXFx0aW1lc1xcbG5yW1xcQV8xXFwsXFxBX247XFxCXVxcdGltZXNcXGxucltcXEFfMVxcLFxcQV9uO1xcQl0iXSxbMiwyLCJcXEFfMVxcbnRpbWVzXFxBX25cXHRpbWVzXFxsbnJbXFxBXzFcXCxcXEFfbjtcXEJdXFx0aW1lc1xcQV8xXFxudGltZXNcXEFfblxcdGltZXNcXGxucltcXEFfMVxcLFxcQV9uO1xcQl0iXSxbMiwxLCJcXEJcXHRpbWVzXFxCIl0sWzEsMCwiXFxldiJdLFszLDQsIlxcZXZcXHRpbWVzXFxldiIsMl0sWzQsMCwiKyIsMl0sWzIsMSwiXFxpZFxcdGltZXMrIl0sWzIsMywiKFxcRGVsdGFcXHRpbWVzXFxpZClcXGV4IiwyXV0=
\begin{tikzcd}
{\A\times\lnr[\A_1\,\A_n;\B]} && \B \\
&& {\B\times\B} \\
{\A_1\ntimes\A_n\times\lnr[\A_1\,\A_n;\B]\times\lnr[\A_1\,\A_n;\B]} && {\A_1\ntimes\A_n\times\lnr[\A_1\,\A_n;\B]\times\A_1\ntimes\A_n\times\lnr[\A_1\,\A_n;\B]}
\arrow["\ev", from=1-1, to=1-3]
\arrow["{+}"', from=2-3, to=1-3]
\arrow["{\id\times+}", from=3-1, to=1-1]
\arrow["{(\Delta\times\id)\ex}"', from=3-1, to=3-3]
\arrow["{\ev\times\ev}"', from=3-3, to=2-3]
\end{tikzcd}
}
\end{equation*}
commutative, where $\Delta\colon M\to M\times M$ is the diagonal map and $\ex\colon M\times M\times N\times N\to M\times N\times M\times N$ denotes the  map that exchanges the internal terms. To check that this morphism is well-defined, we shall prove that $(\Delta\times\id)\ex(\ev\times\ev)+$ is a parametrized multilinear map by computing the following:
\begin{align*}
(\Delta\times\id)\ex(\ev&\times\ev)+\lambda\\
&=~(\Delta\times\id)\ex((\ev\lambda)\times(\ev\lambda))\T+     \tag{Additivity of $\lambda$}\\
&=~\<\lambda_k,z,\lambda_k,z\>\T((\ev\times\ev)+)      \tag{\textbf{[ILH.1']}}\\
&=~(\lambda_k\times z)\T((\Delta\times\id)\ex(\ev\times\ev)+)
\end{align*}
We leave it to the reader to prove that $(\lnr[\A_1\,\A_n;\B],0,+)$ forms a commutative monoid, which in the internal language is equivalent to proving the following equations:
\begin{align*}
&(f(x_1\,x_n)+g(x_1\,x_n))+h(x_1\,x_n)=f(x_1\,x_n)+(g(x_1\,x_n)+h(x_1\,x_n))\\
&f(x_1\,x_n)+g(x_1\,x_n)=g(x_1\,x_n)+f(x_1\,x_n)\\
&f(x_1\,x_n)+0=f(x_1\,x_n)
\end{align*}
where $f,g$, and $h$ are \emph{internal morphisms} in $\lnr[\A_1\,\A_n;\B]$. We shall define the differential projection $\hat p$ as the unique morphism that makes the following diagram
\begin{equation*}
\adjustbox{width=\linewidth}{
% https://q.uiver.app/#q=WzAsNSxbMCwwLCJcXEFfMVxcbnRpbWVzXFxBX25cXHRpbWVzXFxUXFxsbnJbXFxBXzFcXCxcXEFfbjtcXEJdIl0sWzEsMCwiXFxUKFxcQV8xXFxudGltZXNcXEFfblxcdGltZXNcXGxucltcXEFfMVxcLFxcQV9uO1xcQl0pIl0sWzIsMCwiXFxUXFxCIl0sWzIsMSwiXFxCIl0sWzAsMSwiXFxBXzFcXG50aW1lc1xcQV9uXFx0aW1lc1xcbG5yW1xcQV8xXFwsXFxBX247XFxCXSJdLFswLDEsInpfe1xcQV8xXFxudGltZXNcXEFfbn1cXHRpbWVzXFxpZF97XFxUXFxsbnJbXFxBXzFcXCxcXEFfbjtcXEJdfSJdLFsxLDIsIlxcVFxcZXZfe1xcQV8xXFwsXFxBX247XFxCfSJdLFsyLDMsIlxcaGF0IHBfXFxCIl0sWzQsMywiXFxldl97XFxBXzFcXCxcXEFfbjtcXEJ9IiwyXSxbMCw0LCJcXGlkX3tcXEFfMVxcbnRpbWVzXFxBX259XFx0aW1lc1xcaGF0IHBfe1xcbG5yW1xcQV8xXFwsXFxBX247XFxCXX0iLDJdXQ==
\begin{tikzcd}
{\A_1\ntimes\A_n\times\T\lnr[\A_1\,\A_n;\B]} & {\T(\A_1\ntimes\A_n\times\lnr[\A_1\,\A_n;\B])} & {\T\B} \\
{\A_1\ntimes\A_n\times\lnr[\A_1\,\A_n;\B]} && \B
\arrow["{z\times\id}", from=1-1, to=1-2]
\arrow["{\id\times\hat p}"', from=1-1, to=2-1]
\arrow["{\T\ev}", from=1-2, to=1-3]
\arrow["{\hat p}", from=1-3, to=2-3]
\arrow["{\ev}"', from=2-1, to=2-3]
\end{tikzcd}
}
\end{equation*}
commutative. Since $z$ and $\hat p$ are linear, by composition, $(z\times\id)\T\ev\hat p$ is a parametrized multilinear map. Thus, $\hat p$ is well-defined. We leave it to the reader to prove the compatibilities between $\hat p$ and the additive structure of $\lnr[\A_1\,\A_n;\B]$, which is a consequence of the compatibilities between the differential projection $\hat p$ and the additive structure of $\B$. Instead, we focus on proving the universal property that establishes that $\T\lnr[\A_1\,\A_n;\B]$ is isomorphic to $\lnr[\A_1\,\A_n;\B]\times\lnr[\A_1\,\A_n;\B]$ with projections $p$ and $\hat p$. Consider two morphisms $f,g\colon M\to\lnr[\A_1\,\A_n;\B]$. Our goal is to construct a unique morphism $\<f,g\>\colon M\to\T\lnr[\A_1\,\A_n;\B]$ that satisfies the following two equations: $\<f,g\>p=f$ and $\<f,g\>\hat p=g$. We begin by considering the following two morphisms:
\begin{align*}
&\tilde f\colon\A_1\ntimes\A_n\times M\xrightarrow{\id\times f}\A_1\ntimes\A_n\times\lnr[\A_1\,\A_n;\B]\xrightarrow{\ev}\B\\
&\tilde g\colon\A_1\ntimes\A_n\times M\xrightarrow{\id\times g}\A_1\ntimes\A_n\times\lnr[\A_1\,\A_n;\B]\xrightarrow{\ev}\B
\end{align*}
Since $\ev$ is a parametrized multilinear map, it follows that $\tilde f$ and $\tilde g$ are also parametrized multilinear maps. Using that $\B$ is a differential object, we can define a morphism $\<\tilde f,\tilde g\>\colon\A_1\ntimes\A_n\times M\to\T\B$ as the unique morphism that satisfies the following two equations: $\<\tilde f,\tilde g\>p=\tilde f$ and $\<\tilde f,\tilde g\>\hat p=\tilde g$. We shall prove that $\<\tilde f,\tilde g\>$ is a parametrized multilinear map by computing the following:
\begin{align*}
\<\tilde f,\tilde g\>\T\lambda c=\<\tilde f\lambda,\tilde g\lambda\>=(\lambda_k\times z_M)\T(\<\tilde f,\tilde g\>)
\end{align*}
where we used that $\T\lambda c$ becomes $\lambda\times\lambda$ when $\T\B$ is seen as a Cartesian product. Thus, since $\<\tilde f,\tilde g\>$ is a parametrized multilinear map, by \textbf{[ILH.2']}, there exists a unique $\<f,g\>\colon M\to\T\lnr[\A_1\,\A_n;\B]$ that satisfies the following equation:
\begin{align*}
&(z\times\<f,g\>)\T\ev=\<\tilde f,\tilde g\>
\end{align*}
We now prove that $\<f,g\>$ satisfes the desired equations:
\begin{align*}
(\id\times(\<f,g\>p))\ev&=~(z\times\<f,g\>)(p\times p)\ev  \tag{$zp=\id$}\\
&=~(z\times\<f,g\>)\T\ev p     \tag{Naturality of $p$}\\
&=~\<\tilde f,\tilde g\>p                    \tag{Definition of $\<f,g\>$}\\
&=~\tilde f                                   \tag{Definition of $\<\tilde f,\tilde g\>$}\\
&=~(\id\times f)\ev              \tag{Definition of $\tilde f$}
\end{align*}
Thus, by the universal property \textbf{[ILH.2]}, $\<f,g\>p=f$. Now, let us show that $\<f,g\>\hat p=g$:
\begin{align*}
(\id\times(\<f,g\>\hat p))\ev&=~(\id\times\<f,g\>)(\id\times\hat p)\ev\\
&=~(\id\times\<f,g\>)(z\times\id)\T\ev\hat p     \tag{Definition of $\hat p$}\\
&=~(z\times\<f,g\>)\T\ev\hat p\\
&=~\<\tilde f,\tilde g\>\hat p        \tag{Definition of $\<f,g\>$}\\
&=~\tilde g                                   \tag{Definition of $\<\tilde f,\tilde g\>$}\\
&=~(\id\times g)\ev              \tag{Definition of $\tilde g$}
\end{align*}
Thus, using again \textbf{[ILH.2]}, we conclude that $\<f,g\>\hat p=g$. Now, let $h\colon M\to\T\lnr[\A_1\,\A_n;\B]$ satisfy $hp=f$ and $h\hat p=g$. Thus, following the same steps as before, we can easily prove that:
\begin{align*}
&(z\times h)\T\ev p=(\id\times f)\ev=\tilde f    &&(z\times h)\T\ev\hat p=(\id\times g)\ev=\tilde g
\end{align*}
However, there is only one morphism that satisfies these two equations, that is:
\begin{align*}
&(z\times h)\T\ev=\<\tilde f,\tilde g\>=(\<z\times\<f,g\>)\T\ev
\end{align*}
By \textbf{[ILH.2']}, we conclude that $h=\<f,g\>$. Finally, by \textbf{[ILH.1']}, the evaluation map is a parametrized multilinear map. Furthermore, by construction, $(z\times\hat p)\ev=\T\ev\hat p$, which is equivalent to linearity in the last argument. Thus, $\ev$ is multilinear.
\end{proof}

Theorem~\ref{theorem:internal-linear-homs-classify-kleisli-multilinear-morphisms} establishes what class of morphisms is truly classified by an internal multilinear Hom. Before moving on, we shall express condition \textbf{[ILH.2']} only in terms of \textbf{[ILH.2]}, but in a different category. To this end, recall that the tangent bundle functor $\T$ of a tangent category comes with a monad structure, whose unit is the zero vector field $z_M\colon M\to\T M$ and whose multiplication is $\mu_M\colon\T\T M\to\T M$, given by $\<\T p_M,p_{\T M}\>s_M$~\cite[Proposition~3.4]{cockett:tangent-cats}. Furthermore, the canonical flip $c_M\colon\T\T M\to\T\T M$ turns $\T$ into a strong tangent monad, which consists of a monad equipped with a distributive law between the tangent bundle functor and the monad (see~\cite[Definition~19]{cockett:tangent-monads}).

In particular, we shall call a \textbf{colax tangent monad} a monad $S\colon\X\to\X$ equipped with a distributive law $\alpha_M\colon\T SM\to S\T M$, compatible with the monad and the tangent structure; furthermore, we shall assume that the endofunctor $S$ preserves the fundamental tangent limits\footnote{These are the $n$-fold pullback of the projection along itself together with the image of this $n$-fold pullback along $\T^n$, and the pullback diagram that establishes the universal property of the vertical lift.}.

\begin{lemma}
\label{lemma:kleisli-category-of-tangent-bundle-functor}
The Kleisli category of a colax tangent monad $(S,\alpha)$ becomes a tangent category whose tangent bundle functor $\T^\KL$ sends each object $M$ to $\T M$ and each $f\colon M\nto N$ to $\T^\KL f\colon\T M\nto\T N$, whose intepretation is given by:
\begin{align*}
&\SQ{\T^\KL f}\colon\T M\xrightarrow{\T\SQ{f}}\T SN\xrightarrow{\alpha_N}S\T N
\end{align*}
Furthermore, the inclusion functor $(\X,\TT)\to\KL[S,\alpha]$ is a strict tangent morphism that preserves the limits preserved by $S$.
\end{lemma}
\begin{proof}
It is straightforward to lift the structural natural transformations of the tangent structure $\TT$ on the base tangent category to the Kleisli category. For instance, the projection $p^\KL$ is defined by $\SQ{p^\KL}\=p\eta$. The inclusion functor $\X\to\KL[S]$ sends each object $M$ to itself, and each morphism $f\colon M\to N$ to $f^\KL$ defined by $\SQ{f^\KL}\=f\eta$, where $\eta$ is the unit of $S$, and it preserves the limits preserved by the endofunctor $S$ itself. From the functoriality of the inclusion functor, all the equational axioms hold for $\TT^\KL$ to be a tangent structure; from the preservation of the fundamental tangent limits, also the universal properties of the fundamental tangent limits hold. Finally, by construction, the inclusion functor becomes a strict tangent morphism that preserves the limits preserved by $S$.
\end{proof}

\begin{remark}
\label{remark:kleisli-tangent-category}
In~\cite[Proposition~3.5]{lanfranchi:lie-correspondence-tangent-cats}, the author already proved that the \emph{co}Kleisli category of a \emph{lax} tangent \emph{co}monad is a tangent category and that the inclusion functor is a strict tangent morphism. One could see Lemma~\ref{lemma:kleisli-category-of-tangent-bundle-functor} as the dual statement of this proposition. However, one should also keep in mind that the dual argument does not apply directly, since the opposite category of a tangent category fails in general to be a tangent category. In particular, for Lemma~\ref{lemma:kleisli-category-of-tangent-bundle-functor} to hold, we required the monad $S$ to preserve the fundamental tangent limits, a condition that is not required for lax tangent comonads.
\end{remark}

Since $(\T,c)$ is a strong tangent monad, $(\T,c)$ becomes a colax tangent monad; so, in particular, its Kleisli category is a tangent category. Furthermore, as proved in~\cite[Corollary~5.31]{lanfranchi:phd-thesis}, not only $(\T,c)$ is a strong tangent monad, but so is every $(\T^n,c^n)$, where $c^n$ is given by $c^n\=\T^nc\T^{n-1}c_\T\dots\T c_{\T^{n-1}}c_{\T^n}$.

\begin{proposition}
\label{proposition:kleisli-tangent-category-of-tangent-bundle-functor}
Given a (Cartesian) tangent category $(\X,\TT)$, for each $n$, the Kleisli category $\KL[\T^n,c^n]$ of $(\T^n,c^n)$ is a (Cartesian) tangent category and the inclusion functor $(\X,\TT\to\KL[\T^n,c^n]$ is a strict tangent morphism (that preserves Cartesian products). In particular, it preserves differential objects and multilinear maps.
\end{proposition}

We can now express condition \textbf{[ILH.2']} in a clearer form.

\begin{corollary}
\label{corollary:internal-linear-homs-classify-kleisli-multilinear-morphisms}
Given a list $\A_1\,\A_n,\B$ of differential objects, consider an object $\lnr[\A_1\,\A_n;\B]$ equipped with a map $\ev\colon\A_1\ntimes\A_n\times\lnr[\A_1\,\A_n;\B]\to\B$. Then $\lnr[\A_1\,\A_n;\B]$ is an internal multilinear Hom with evaluation map $\ev$ if and only if the following conditions are satisfied:
\begin{description}
\item[ILH.1'] $\ev$ is a $\lnr[\A_1\,\A_n;\B]$-parametrized multilinear map;

\item[ILH.2''] The morphism $\ev^\KL$ defined by $\SQ{\ev^\KL}=\ev z$, classifies multilinear maps in the Kleisli tangent category $\KL[\T,c]$, in that for every $M$-parametrized multilinear map $f\colon\A_1\ntimes\A_n\times M\nto\B$, there exists a unique morphism $\hat f\colon M\nto\lnr[\A_1\,\A_n;\B]$ that makes the following diagram
\begin{equation*}
% https://q.uiver.app/#q=WzAsMyxbMCwwLCJcXEFfMVxcbnRpbWVzXFxBX25cXHRpbWVzXFxsbnJbXFxBXzFcXCxcXEFfbjtcXEJdIl0sWzEsMCwiXFxCIl0sWzAsMSwiXFxBXzFcXG50aW1lc1xcQV9uXFx0aW1lcyBNIl0sWzAsMSwiXFxldl5cXEtMIiwwLHsic3R5bGUiOnsiYm9keSI6eyJuYW1lIjoiYmFycmVkIn19fV0sWzIsMSwiZiIsMix7InN0eWxlIjp7ImJvZHkiOnsibmFtZSI6ImJhcnJlZCJ9fX1dLFsyLDAsIlxcaWRcXHRpbWVzXFxoYXQgZiIsMCx7InN0eWxlIjp7ImJvZHkiOnsibmFtZSI6ImJhcnJlZCJ9fX1dXQ==
\begin{tikzcd}
{\A_1\ntimes\A_n\times\lnr[\A_1\,\A_n;\B]} & \B \\
{\A_1\ntimes\A_n\times M}
\arrow["{\ev^\KL}"{inner sep=.8ex}, "\shortmid"{marking}, from=1-1, to=1-2]
\arrow["{\id\times\hat f}"{inner sep=.8ex}, "\shortmid"{marking}, from=2-1, to=1-1]
\arrow["f"'{inner sep=.8ex}, "\shortmid"{marking}, from=2-1, to=1-2]
\end{tikzcd}
\end{equation*}
commutative in $\KL[\T,c]$.
\end{description}
\end{corollary}
\begin{proof}
To prove this result, we shall simply rewrite condition \textbf{[ILH.2']} as \textbf{[ILH.2'']}. To this end, consider an $M$-parametrized multilinear map $f\colon\A_1\ntimes\A_n\times M\to\T\B$. This is clearly a morphism in the Kleisli category $\KL[\T,c]$. On the other hand, $\ev^\KL$ is defined by $\SQ{\ev^\KL}=\ev z$. Thus, by Kleisli composition, which we shall denote by $;_\KL$, for a given morphism $\hat f\colon M\to\T\lnr[\A_1\,\A_n;\B]$:
\begin{align*}
\SQ{(\id\times\hat f);_\KL\ev^\KL}&=~(\SQ{\id}\times\hat f)\T\SQ{\ev^\KL}\mu    \tag{Kleisli composition and products in the Kleisli}\\
&=~(z\times\hat f)\T\ev\T z\mu                                                  \tag{Identity in the Kleisli category and $\SQ{\ev^\KL}=\ev z$}\\
&=~(z\times\hat f)\T\ev                                                         \tag{Unitality: $\T z\mu=\id$}
\end{align*}
Therefore, $(\id\times\hat f);_\KL\ev^\KL=f$ in $\KL[\T,c]$ if and only if $(z\times\hat f)\T\ev=f$ in the base tangent category.
\end{proof}

We shall refer to the morphisms in the Kleisli category as \textbf{Kleisli morphisms}. We shall also refer to the $M$-parametrized multilinear morphisms in the Kleisli tangent category $\KL[\T^n,c^n]$ as \textbf{Kleisli $M$-parametrized multilinear morphisms} of order $n$.



%__________________________________________________________________________
\subsection{The algebraic structures of Linear Endomorphism Classifiers}
\label{subsection:ring-lie-algebra-lec}
For the story of this paper, instead of considering every internal multilinear Hom $\lnr[\A_1\,\A_n;\B]$, we shall focus on the special case $\lnr[\A,\A]$.

\begin{definition}
\label{definition:linear-endomorphism-classifier}
A differential object $\A$ admits the \textbf{Linear Endomorphism Classifier} provided that the internal multilinear Hom $\lnr_\A\=\lnr[\A,\A]$ exists. we will refer to $\lnr_\A$ as the \textbf{Linear Endomorphism Classifier} (\textbf{LEC}) of $\A$.
\end{definition}

In this section, we study the algebraic structures of LEC objects. We start with the most obvious one: every LEC is a rig object\footnote{Rigs are also known in the literature as semirings. Here, we shall not assume commutativity of the multiplication.}. Recall that a \textbf{differential rig object} consists of a differential object $\R$ equipped with a unit $\e\colon\*\to\R$ and a bilinear multiplication map $\m\colon\R\times\R\to\R$ which is associative and unital (see~\cite[Definition~5.24]{cockett:differential-equations}). We notice that such an object is equivalent to a monoid in the multicategory $\DO_\multi[\X,\TT]$.

\begin{proposition}
\label{proposition:rig-structure-on-lec}
The LEC object $\lnr_\A$ of a differential object $A$ is a differential rig object whose unit and bilinear multiplication $\m\colon\lnr_\A\times\lnr_\A\to\lnr_\A$ are respectively defined by:
\begin{equation*}
% https://q.uiver.app/#q=WzAsMyxbMCwwLCJcXEFcXHRpbWVzXFxsbnJfXFxBIl0sWzEsMCwiXFxBIl0sWzAsMSwiXFxBXFx0aW1lc1xcKiJdLFsyLDEsIlxccGlfMSIsMl0sWzAsMSwiXFxldl9cXEEiXSxbMiwwLCJcXGlkX1xcQVxcdGltZXNcXGUiXV0=
\begin{tikzcd}
{\A\times\lnr_\A} & \A \\
{\A\times\*}
\arrow["{\ev_\A}", from=1-1, to=1-2]
\arrow["{\id_\A\times\e}", from=2-1, to=1-1]
\arrow["{\pi_1}"', from=2-1, to=1-2]
\end{tikzcd}\qquad
% https://q.uiver.app/#q=WzAsNCxbMCwwLCJcXEFcXHRpbWVzXFxsbnJfXFxBIl0sWzAsMSwiXFxBXFx0aW1lc1xcbG5yX1xcQVxcdGltZXNcXGxucl9cXEEiXSxbMSwxLCJcXEFcXHRpbWVzXFxsbnJfXFxBIl0sWzEsMCwiXFxBIl0sWzAsMywiXFxldl9cXEEiXSxbMiwzLCJcXGV2X1xcQSIsMl0sWzEsMiwiXFxldl9cXEFcXHRpbWVzXFxpZF97XFxsbnJfXFxBfSIsMl0sWzEsMCwiXFxpZF9cXEFcXHRpbWVzXFxtIl1d
\begin{tikzcd}
{\A\times\lnr_\A} & \A \\
{\A\times\lnr_\A\times\lnr_\A} & {\A\times\lnr_\A}
\arrow["{\ev_\A}", from=1-1, to=1-2]
\arrow["{\id_\A\times\m}", from=2-1, to=1-1]
\arrow["{\ev_\A\times\id_{\lnr_\A}}"', from=2-1, to=2-2]
\arrow["{\ev_\A}"', from=2-2, to=1-2]
\end{tikzcd}
\end{equation*}
\end{proposition}
\begin{proof}
By Proposition~\ref{proposition:exponential-objects}, $\lnr_\A=\lnr[\A,\A]$ is an exponential object in the multicategory $\DO_\multi[\X,\TT]$. Thus, it becomes a monoid object in the multicategory whose multiplication $\A^\A,\A^\A\to\A^\A$ is the uncurry of $(\ev_\A\times\A)\ev_\A\colon\A\times\A^\A\times\A^\A\to\A$ and whose unit $\text{\textvisiblespace}\to\A^\A$ is the uncurry of the identity $\id\colon\A\to\A$.
\end{proof}

One can turn any ring $R$ into a Lie algebra by considering the Lie bracket given by the commutator $[r,s]\=sr-rs$. Unfortunately, in order to define the commutator, one needs negatives, which are not assumed to exist in a rig. To solve this issue, we shall consider the maximal subring of a rig object.

\begin{definition}
\label{definition:maximal-differential-subring}
In a Cartesian tangent category, the \textbf{maximal differential subring} of a differential rig $\R$ consists of a differential object $\R^-$ together with a linear morphism $\varphi\colon\R^-\to\R\times\R$ that turns the following diagram
\begin{equation*}
% https://q.uiver.app/#q=WzAsNCxbMSwwLCJcXFJcXHRpbWVzXFxSIl0sWzMsMCwiXFxSIl0sWzIsMCwiXFwqIl0sWzAsMCwiXFxSXi0iXSxbMCwxLCIrIiwwLHsib2Zmc2V0IjotMn1dLFswLDIsIiIsMix7Im9mZnNldCI6Mn1dLFsyLDEsIjAiLDIseyJvZmZzZXQiOjJ9XSxbMywwLCJcXHZhcnBoaSIsMCx7InN0eWxlIjp7ImJvZHkiOnsibmFtZSI6ImRhc2hlZCJ9fX1dXQ==
\begin{tikzcd}
{\R^-} & {\R\times\R} & {\*} & \R
\arrow["\varphi", dashed, from=1-1, to=1-2]
\arrow[shift right=2, from=1-2, to=1-3]
\arrow["{+}", shift left=2, from=1-2, to=1-4]
\arrow["0"', shift right=2, from=1-3, to=1-4]
\end{tikzcd}
\end{equation*}
into an equalizer.
\end{definition}

The maximal differential ring of a differential rig becomes a differential ring that turns the inclusion morphism $\rho\colon\R^-\to\R$ into a linear rig homomorphism.

\begin{lemma}
\label{lemma:maximal-differential-subring}
The maximal differential ring $\R^-$ of a differential rig $\R$ comes with a unique differential ring structure that makes $\rho\=\varphi\pi_1\colon\R^-\to\R$ into a linear rig homomorphism.
\end{lemma}
\begin{proof}
This result applies more generally to rig objects and their maximal subring objects and is an easy exercise of categorical algebra. Concretely, to construct the negation morphism $-\colon\R^-\to\R^-$ one considers the morphism $\varphi\pi_2$, where $\varphi\colon\R^-\to\R\times\R$ denotes the equalizer morphism. Since $\varphi+=!0$, it follows that $\<!0,\varphi\pi_2\>+=\<!0\varphi\pi_1,\varphi\pi_2\>+=\<!0,\id\>\varphi+=!0$. Therefore, by the universal property of the equalizer, there exists a unique morphism $-\colon\R^-\to\R^-$ such that $-\varphi=\<!0,\varphi\pi_2\>$. One can show that $-$ defines the negation. We leave it to the reader to show that for each rig object $R$, the equalizer $R^-$ of $+$ and $0$ is in fact a ring object. Furthermore, it is also straightforward to see that the inclusion map $\rho$ becomes a rig homomorphism and that, by the universal property, the ring structure of $\R^-$ is uniquely determined by $\rho$ being a rig homomorphism. Finally, to conclude, we need to prove that $\R^-$ is a \emph{differential} ring. To this end, we shall use that $\R^-$ is a differential object and that $\varphi\colon\R^-\to\R\times\R$ is linear. Since $\varphi$ is linear, it follows that $\rho=\varphi\pi_1\colon\R^-\to\R$ is linear as well. Let $\m^-$ denote the multiplication of $\R^-$. Using that $\rho$ is a linear rig homomorphism, and that $\m$ is bilinear, it is not hard to prove that:
\begin{align*}
&(\id\times z)\T\m^-\hat p^-\rho=(\hat p^-\times\id)\m^-\rho &&(z\times\id)\T\m^-\hat p^-\rho=(\id\times\hat p^-)\m^-\rho
\end{align*}
Furthermore, using that $\varphi=\<\rho,-\rho\>$ since $-\varphi=\<!0,\varphi\pi_2\>$, it follows that $\rho$ is monic and therefore, $\m^-$ becomes bilinear.
\end{proof}

Every differential ring comes with a differential Lie algebra structure, induced by the commutator. For starters, let us recall the notion of differential Lie algebra.

\begin{definition}[{\cite[Definition~5.1]{lanfranchi:lie-correspondence-tangent-cats}}]
\label{definition:differential-lie-algebra}
In a Cartesian tangent category, a \textbf{differential Lie algebra} consists of a differential object $\g$ equipped with a morphism $[,]\colon\g\times\g\to\g$ called the \textbf{Lie bracket} of $\g$, satisfying the following conditions:
\begin{description}
\item[DLIE.1] $(\g,0,+,[,])$ is a Lie algebra, in that\footnote{For our convenience, we use pointwise notation to express antisymmetry and the Jacobian identity, as if $x,y$ and $z$ were internal variables of $\g$; however, it should be clear that these axioms can also be fully expressed diagrammatically.}:
\begin{align*}
&[x,y]+[y,x]=0              &&[[x,y],z]+[[z,x],y]+[[y,z],x]=0
\end{align*}

\item[DLIE.2] $[,]\colon\g\times\g\to\g$ is bilinear.
\end{description}
\end{definition}

\begin{proposition}
\label{proposition:LEC-lie-algebra}
The commutator $[,]$ of a differential ring $\R$ equips $\R$ with a differential Lie algebra structure. In particular, the maximal differential subring of a differential rig is a differential Lie algebra.
\end{proposition}
\begin{proof}
To prove that $[,]$ satisfies the antisymmetry identity, we shall use an internal language proof:
\begin{align*}
&[x,y]+[y,x]=\m^\op(x,y)-\m(x,y)+\m^\op(y,x)-\m(y,x)=\\
&\qquad=\m(y,x)-\m(x,y)+\m(x,y)-\m(y,x)=0
\end{align*}
Similarly, using the associativity of $\m$, one can prove the Jacobi identity with a similar technique. We leave the details to the reader. Finally, since $\m$ is bilinear, it is not hard to see that the linear combination $\m^\op-\m$ is also bilinear, proving that $\lnr_\A$ is in fact a differential Lie algebra.
\end{proof}

In particular, the maximal differential subring $\lnr_\A^-$ of the LEC object $\lnr_\A$ of a differential object $\A$ becomes a differential Lie algebra with Lie bracket given by the commutator.



%__________________________________________________________________________
%__________________________________________________________________________